\documentclass[11pt,a4paper]{article}

% ---------- Codifica e lingua ----------
\usepackage[T1]{fontenc}
\usepackage[utf8]{inputenc}
\usepackage[main=italian,provide=*]{babel}

% ---------- Matematica ----------
\usepackage{amsmath,amssymb}
\usepackage{mathtools}

% ---------- Impaginazione ----------
\usepackage[a4paper,margin=2.7cm]{geometry}
\usepackage{caption}
\usepackage{booktabs}
\usepackage{enumitem}
\usepackage{placeins}
\usepackage[colorlinks=true,linkcolor=black,citecolor=black,urlcolor=blue!55!black]{hyperref}

\newcommand{\ket}[1]{\lvert #1 \rangle}
\newcommand{\bra}[1]{\langle #1 \rvert}

\title{Il trimero ad anello\\
\large Panoramica della pipeline (vista a 2000 piedi)}
\author{Samuele \\ Relatore: Prof.\ Alessandro Chiesa}
\date{}

\begin{document}
\maketitle

\begin{abstract}
\noindent Questo documento è una mappa, non una derivazione: serve a
vedere in un colpo d'occhio come è strutturata l'intera pipeline sul
trimero ad anello --- sistema chiuso e poi sistema aperto --- prima di
entrare nel dettaglio di ciascuno stadio (ciascuno ha un proprio
documento dedicato). Copre entrambe le parti in un'unica narrazione
continua: Parte 1 (Sez.~\ref{sec:parte1}), il sistema chiuso, stato puro
$\ket\psi$ che evolve unitariamente; Parte 2 (Sez.~\ref{sec:parte2}), lo
stesso sistema reso aperto dal rumore, operatore densità $\rho$. È il
mirror esatto della pipeline già completata sul dimero: stesso schema in
quattro-più-cinque stadi, stesso modello di rumore, stessa terminologia.
Cambia il sistema --- da 2 a 3 spin, geometria a triangolo isoscele
anziché un solo legame, ansatz $W$-2q.6 al posto di PMA-2q$\cdot$3 ---
non la struttura della pipeline. Scope e canali di rumore restano quelli
confermati dal relatore; topologia confermata (6 settembre 2026):
l'\textbf{anello}, non la catena.
\end{abstract}

\tableofcontents

\section{Parte 1 --- Il sistema chiuso}
\label{sec:parte1}

\subsection{Da dove partiamo: il trimero isoscele frustrato}

Il dimero (un solo legame, due spin) non poteva porre alcune domande per
mancanza di struttura: non c'è frustrazione con un solo legame, non c'è
scelta di \emph{forma} del termine DM con un solo legame, non c'è
differenza fra ``il punto giusto per il VQE'' e ``il punto giusto per
mostrare una dinamica non banale''. Il trimero ad anello --- tre spin
$1/2$, geometria isoscele: legame $J$ fra i siti $(1,2)$, legami $J'$ fra
$(2,3)$ e $(3,1)$ --- introduce tutte e tre. Ogni spin $1/2$ si mappa
direttamente su un qubit, come sul dimero: nessun costo di
rappresentazione aggiuntivo nel passare da due a tre siti. Con
$J'\neq J$ il ciclo è genuinamente frustrato (a differenza del caso
equilatero $J'=J$, dove la frustrazione è presente ma degenere in modo
diverso) --- è questo legame in più, assente nella catena aperta, a
rendere il sistema geometricamente frustrato per costruzione, non per
una scelta di parametri ad hoc.

\subsection{Il termine DM: due opzioni, una scelta}
\label{subsec:dm-parte1}

A differenza del dimero, dove il DM ha una forma unica (un solo legame,
nessuna scelta), con tre legami non c'è un solo modo di accendere un
termine DM --- e non è una scelta libera: l'intera teoria del trimero
isoscele assume l'invarianza per scambio dei siti $1,2$. Verificato
sistematicamente (non assunto): costruita esplicitamente la matrice di
permutazione $P_{12}$ e controllato il commutatore $\|[P_{12},H_\text{DM}]\|$
per tutte le 12 combinazioni di segno di $(D_{12},D_{23},D_{31})$
proporzionali a $(J,J',J')$ --- fra tutte, \textbf{una sola} coppia
rispetta la simmetria esattamente (nessun DM sul legame di base, segno
opposto sui due legami laterali): è l'\textbf{Opzione A}. Il candidato
proposto dal relatore (tutti e tre i legami, stesso segno) --- l'
\textbf{Opzione B} --- non solo rompe quella simmetria: è, fra tutte le
combinazioni testate, quella che la rompe di più.

La ragione per cui l'Opzione A lascia l'incrocio fondamentale/primo
eccitato genuinamente chiuso (mentre l'Opzione B lo apre) non è una
coincidenza numerica: è la regola di non-incrocio di von
Neumann--Wigner applicata a $S_{12}^2$, la simmetria residua che
l'Opzione A rispetta e l'Opzione B rompe. Per il resto della pipeline
--- VQE, dinamica, correlatori --- si usa sempre l'Opzione B: è l'unica
che riproduce, sul trimero, lo stesso ruolo che il DM ha già sul dimero
(rendere unico lo stato fondamentale al punto critico, condizione
necessaria per un VQE ben posto).

\subsection{Gli stadi di Parte 1}
\label{sec:stadi-parte1}

\begin{table}[htbp]
\centering
\begin{tabular}{@{}cll@{}}
\toprule
\# & Stadio & Cosa introduce di nuovo rispetto al dimero \\
\midrule
1 & Il sistema & frustrazione geometrica, due opzioni DM \\
2 & Il VQE & due famiglie di ansatz a confronto (RBS, $W$) \\
3 & La dinamica (Trotter) & tre livelli annidati, non due \\
4 & I correlatori & 81 combinazioni, non 36 \\
\bottomrule
\end{tabular}
\caption{I quattro stadi di Parte 1. Ciascuno ha un proprio documento
dedicato nella serie narrativa (Documenti 0--4).}
\end{table}
\FloatBarrier

\subsubsection{Il sistema}

Diagonalizzazione esatta della matrice $8\times8$, decomposizione di
Kambe in tre multipletti ($S_{12}=0$, blocco C; $S_{12}=1$, blocchi A e
B, rispettivamente $S=3/2$ e $S=1/2$) --- stessa tecnica del dimero,
estesa a una famiglia in più. Il risultato che conta per il resto del
lavoro: l'incrocio fondamentale/primo eccitato (fra blocco C e blocco A)
è protetto dalla conservazione di $S_{12}^2$ finché nessun termine la
rompe --- da cui discende direttamente la Sez.~\ref{subsec:dm-parte1}.

\subsubsection{Il VQE}

Due famiglie di ansatz a confronto, stessa struttura esterna (sei
parametri: tre sui legami, tre rotazioni $R_y$ finali indipendenti):
RBS-2q.6 (rotazione di Givens reale, genuinamente $M$-conservante) e
$W$-2q.6 (due CNOT attorno a una $R_y$, lo stesso blocco che sul dimero
si chiama ``ansatz base''). Senza DM, equivalenti ($\mathcal F=1$ per
entrambe). Con DM acceso, al punto di lavoro Scenario A ($b/J=b_c=2.4$,
$D/J=0.15$, Opzione B): RBS-2q.6 ha un tetto strutturale vero
($\mathcal F=0.9999110739$, non si ripara aggiungendo giri); $W$-2q.6
raggiunge l'esatto ($\mathcal F=1$ a precisione macchina), con meno gate
--- l'opposto di quanto succede sul dimero, dove RBS è la scelta
economica e $W$ (la famiglia originale di Crippa et al.) è più costosa
senza vantaggi. Verificato inoltre che $W$-2q.6 non ha bisogno di rami
di preparazione diversi a seconda di $b/J$ (13 punti testati, incluso
l'incrocio stesso, $\mathcal F=1$ ovunque) --- a differenza dell'ansatz
minimo del dimero.

\subsubsection{La dinamica (Trotter)}

Struttura a \textbf{tre} livelli annidati, non due come sul dimero:
campo e scambio isotropo fattorizzano esattamente fra loro ($[S_z^\text{tot},H_\text{ex}]=0$),
ma lo scambio stesso richiede un Trotter interno sui tre legami
condivisi (i legami $(1,2)$ e $(2,3)$ condividono il qubit del sito 2) ---
e lo stesso vale per il termine DM. Convergenza molto più lenta del
dimero a parità di soglia di accuratezza: servono centinaia di passi,
non decine, per via dei due livelli interni aggiuntivi. Due punti di
lavoro distinti (a differenza del dimero, dove uno solo serve entrambi
gli scopi): Scenario A ($b=b_c=2.4$, $D=0.15$, quello del VQE) e il
punto $R_0$ ($b=0.05$, $D=1.93$, \emph{dichiarato provvisorio}), scelto
apposta per una dinamica non monocromatica, dove l'errore di Trotter è
osservabile.

\subsubsection{I correlatori}

Stesso schema Hadamard test del dimero, registro più grande (tre qubit
fisici più l'ancilla, quattro in tutto; ancilla = qubit 3, non 0 come
sul dimero). Con tre siti e tre componenti le combinazioni non sono più
36 ma \textbf{81}. Stesso argomento del dimero per gli zeri all'istante
iniziale (32/81, stato fondamentale reale); ma solo 4/81
\textbf{strutturalmente} nulli per ogni $t$, tutti sullo stesso sito
(sito 3), non ancora spiegati in forma chiusa. Chiuso anche il
confronto con la preparazione reale (VQE $W$-2q.6 contro ampiezze
esatte): scarto medio $2.9\times10^{-7}$, nessuna correlazione con
l'ampiezza --- qui l'ansatz è quasi esatto, non c'è la stessa storia
ricco/monocromatico che il dimero può raccontare con un ansatz
genuinamente imperfetto.

\subsection{Cosa resta fuori (Parte 1)}

Il sistema come sistema aperto (rumore di gate e di lettura) è trattato
in un pacchetto separato (Parte 2, Sez.~\ref{sec:parte2}). La topologia
\textbf{catena aperta} (teoria esatta già disponibile) non è ripresa in
questa serie --- decisione motivata e confermata dal relatore, non persa:
resta un filone possibile per il seguito. Le reti neurali (NQS) sono
fuori scope fin dall'inizio.

\clearpage
\section{Parte 2 --- Il sistema aperto (rumore)}
\label{sec:parte2}

\subsection{Da dove partiamo}

La parte a sistema chiuso è completa sul trimero ad anello
(Sez.~\ref{sec:parte1}): teoria esatta (decomposizione di Kambe, termine
DM Opzione B), VQE senza e con DM (ansatz $W$-2q.6), evoluzione
temporale $e^{-iHt}$ via Suzuki--Trotter, e correlazioni dinamiche
$C_{ij}^{\alpha\beta}(t)$ tramite lo stesso circuito di tipo Hadamard
test usato sul dimero. Tutto lavora su uno stato \emph{puro} $\ket{\psi}$
a 3 qubit che evolve unitariamente: non c'è interazione con l'ambiente,
non c'è rumore.

Questa pipeline ripercorre esattamente la stessa catena di stadi già
fatta sul dimero --- preparazione dello stato fondamentale, evoluzione
temporale, misura dei correlatori --- ma la accompagna con le stesse due
sorgenti di rumore realistiche già calibrate per il dimero. Non si
introduce fisica nuova nell'Hamiltoniana né un nuovo modello di rumore:
si applica lo stesso modello, già validato, a un circuito a 3 qubit
invece che 2.

\subsection{Cosa cambia: da \texorpdfstring{$\ket{\psi}$}{|psi>} a \texorpdfstring{$\rho$}{rho}}

Come per il dimero, il rumore accoppia il sistema a un ambiente che non
osserviamo: uno stato puro non basta più, serve l'operatore densità
$\rho$ e un'evoluzione non unitaria pura. In Qiskit questo corrisponde,
esattamente come per il dimero, a usare
\texttt{AerSimulator(method=\allowbreak"density\_matrix")} invece del
simulatore statevector --- qui su un registro a 3 qubit.

\subsection{Gli stadi della pipeline}
\label{sec:stadi}

\begin{table}[htbp]
\centering
\begin{tabular}{@{}cll@{}}
\toprule
\# & Stadio & Cosa cambia rispetto al sistema chiuso \\
\midrule
1 & Modello di rumore & wrapper sul dimero, stessa logica \\
2 & VQE con termine DM, rumoroso & ansatz $W$-2q.6, simulatore a $\rho$ \\
3 & Quantum simulation (Trotter), rumorosa & stesso circuito, compilato per passo \\
4 & Correlazioni dinamiche, rumorose & stesso Hadamard test, readout incluso \\
5 & Scan sui parametri di rumore & nuovo: nessun analogo nel sistema chiuso \\
\bottomrule
\end{tabular}
\caption{I cinque stadi della pipeline, mirror esatto di quelli del
dimero. Il primo stadio non richiede alcun adattamento della
\emph{fisica} del modello: il \texttt{NoiseModel} non dipende dal
numero di qubit per costruzione (verificato, non assunto ---
Sez.~\ref{subsec:modello}). Un nuovo file esiste comunque
(\texttt{noise\_model\_\allowbreak trimero\_anello.py}), ma è un
wrapper sottile, non una reimplementazione.}
\end{table}
\FloatBarrier

\subsubsection{Modello di rumore}
\label{subsec:modello}

Stessi due canali confermati dal relatore per il dimero (errore di gate
depolarizzante 1q/2q, errore di lettura simmetrico), stessa calibrazione
reale (\texttt{ibm\_torino}, arXiv:2504.15187). La logica del
\texttt{NoiseModel} non viene reimplementata: costruisce l'errore con
\texttt{add\_all\_qubit\_\allowbreak quantum\_error}/\texttt{add\_all\_qubit\_\allowbreak readout\_error},
operazioni intrinsecamente indipendenti dal numero di qubit del circuito
su cui vengono applicate. \texttt{noise\_model\_\allowbreak trimero\_anello.py} è un
\textbf{wrapper sottile} che importa e ri-espone questa logica da
\texttt{noise\_model\_dimero.py} (dipendenza dichiarata esplicitamente
nel docstring, non nascosta), così che il trimero abbia un punto
d'ingresso con lo stesso nome usato in ogni altra parte del progetto,
senza duplicare codice da tenere sincronizzato. Verificato empiricamente
sul circuito VQE+DM a 3 qubit, non solo argomentato dalla forma del
codice. Documento:
\texttt{risultati\_\allowbreak modello\_rumore\_\allowbreak trimero\_anello.tex}.

\subsubsection{VQE con termine DM, sotto rumore}

Si ripete la stima del ground state con l'ansatz $W$-2q.6 già validato a
sistema chiuso (fedeltà $\approx1$ nel caso ideale, punto di lavoro
$J=1,\,J'=0.4,\,b=b_c=2.4,\,D=0.15$), eseguito su simulatore a operatore
densità con il \texttt{NoiseModel} agganciato. Documento:
\texttt{risultati\_vqe\_dm\_\allowbreak rumoroso\_trimero\_anello.tex}.

\subsubsection{Quantum simulation (Trotter), sotto rumore}
\label{subsec:trotter}

Stesso circuito di Suzuki--Trotter del sistema chiuso, stessa insidia
già affrontata sul dimero: transpilare l'\emph{intero} circuito a un
livello di ottimizzazione alto collasserebbe gli $N$ passi identici in
un costo costante, cancellando la dipendenza da $N$ che serve allo scan.
A differenza del dimero, qui la strategia di transpilazione a blocchi
va scritta da zero: né
\texttt{circuito\_\allowbreak correlazioni\_\allowbreak trimero\_anello.py}
né \texttt{trotter\_trimero\_anello.py} la contengono già (verificato,
non assunto). Soluzione: transpilare una volta il singolo passo/blocco
al suo costo minimo, poi ripetere il blocco già compilato $N$ volte.

\textbf{Due scenari trattati}, non uno solo: Parte 1 dell'anello usa
punti di lavoro diversi per il VQE ($b=b_c=2.4,\,D=0.15$) e per la
dimostrazione Trotter ($R_0$: $b=0.05,\,D=1.93$, provvisorio, da
$\ket{000}$) --- a differenza del dimero, dove un unico punto serve
entrambi. Lo Stadio 3 copre entrambi: scenario A (preparazione VQE + $N$
passi) mirror esatto del dimero, scenario B (nessuna preparazione, da
$\ket{000}$) mirror della dimostrazione di Parte 1, reso rumoroso.
Documento:
\texttt{risultati\_\allowbreak trotter\_\allowbreak rumoroso\_trimero\_anello.tex}.

\subsubsection{Correlazioni dinamiche, sotto rumore}

Stesso circuito Hadamard test del sistema chiuso (ancilla + 3 qubit di
registro), ora con rumore su preparazione, blocco $U(t)$ ripetuto e
misura sull'ancilla (readout incluso). Documento:
\texttt{risultati\_correlatori\_\allowbreak rumorosi\_trimero\_anello.tex}.

\subsubsection{Scan sui parametri di rumore}

Come per il dimero: si ripete il conto dello stadio precedente al
variare di $(\varepsilon_{1q},\varepsilon_{2q},p_\text{readout})$, per
vedere come si sposta il numero di passi ottimale $N^*$. Documento:
\texttt{risultati\_scan\_\allowbreak parametri\_rumore\_trimero\_anello.tex}.

\subsection{Le due estensioni facoltative}
\label{sec:estensioni}

Oltre ai cinque stadi principali, due estensioni non prioritarie ma
completate e verificate --- non chiudono la pipeline (che resta chiusa
sui cinque stadi, Sez.~\ref{sec:stadi}), la arricchiscono.

\paragraph{VQE noise-aware.} Invece di riusare i parametri ideali
(Stadio 2), si riottimizza l'energia \emph{direttamente sotto rumore} ---
come si farebbe su hardware vero, dove non esiste un riferimento ideale
per calibrare. A differenza del dimero, dove questa estensione non
trovava alcun vantaggio a riottimizzare, qui \textbf{sì}: $\Delta
E\approx-4.16\times10^{-3}$, $\Delta F\approx-1.58\times10^{-2}$ (a
struttura di circuito fissa). Trovato per strada un artefatto del
transpilatore (non un effetto fisico): a valori speciali di $\theta$
(vicini a $0$ o $\pi$) il transpilatore può ridurre il conteggio CNOT
($6\to4$) --- isolato esplicitamente lavorando a struttura fissa
(livello di ottimizzazione 0) come metodologia principale. $N^*$ a
valle (Trotter e correlatore) resta invariato nonostante lo spostamento
di energia e fedeltà. Documento:
\texttt{risultati\_vqe\_\allowbreak noise\_aware\_\allowbreak trimero\_anello\_definitivo.tex}.

\paragraph{Readout asimmetrico.} Il caso simmetrico ($p_{01}=p_{10}$)
usato negli Stadi 1--5 è una semplificazione: su hardware reale il
rilassamento $T_1$ durante la misura rende tipicamente $p_{10}>p_{01}$.
Verificato con readout genuinamente campionato a shot finiti (non la
formula analitica --- vedi nota sotto): $N^*=3$ invariante su un
rapporto $p_{10}/p_{01}$ da $1\times$ a $50\times$, ben oltre il range
realistico. A differenza dell'estensione precedente, qui nessuna
sorpresa: stesso esito qualitativo del dimero. Documento:
\texttt{risultati\_readout\_\allowbreak asimmetrico\_\allowbreak trimero\_anello\_definitivo.tex}.

\paragraph{Nota sul metodo di lettura, valida per tutta la pipeline.} Gli
Stadi 1--5 (e la prima verifica dell'estensione readout asimmetrico)
applicano l'errore di lettura \textbf{analiticamente}:
$\langle Z\rangle_\text{readout}=(1-2p)\langle Z\rangle_\text{ideale}$.
Verificato successivamente (readout genuinamente campionato, shot
finiti via \texttt{AerSimulator} con \texttt{ReadoutError} vero
agganciato al \texttt{NoiseModel}) che le due strade coincidono entro
l'errore statistico atteso --- non assunto, controllato esplicitamente
prima di usare i shot finiti come riferimento per l'estensione
asimmetrica (dove la formula $(1-2p)$ non si applica più così com'è).

\subsection{Cosa resta fuori (Parte 2)}

Stesso scope del dimero, per coerenza dichiarata esplicitamente:
$T_1/T_2$ esclusi per decisione del relatore; reti neurali (NQS) fuori
scope. In più, per il trimero: la topologia \textbf{catena aperta}
resta un'estensione possibile ma non prioritaria (decisione motivata e
confermata dal relatore, 6 settembre 2026) --- il risultato più
originale della catena (alternanza dei legami che cancella l'errore
chirale di Trotter) resta sul tavolo, non perso.

\subsection{Il filo conduttore verificato in ogni stadio}

Stessa proprietà usata nella pipeline del dimero, da riverificare qui
(non assumere dall'analogia): il readout simmetrico è un fattore
\textbf{moltiplicativo} sull'osservabile misurato, mai additivo --- da
cui, se confermato, discenderebbe che anche qui $N^*$ non dipende da
$p_\text{readout}$. Sul dimero questo era un risultato analitico
generale (non specifico a $N=2$); l'aspettativa è che si trasferisca
inalterato, ma resta da verificare esplicitamente allo stadio 5, non da
dare per scontato solo perché la dimostrazione originale non usava
dettagli specifici del numero di spin.

\begin{quote}
\textit{Confermato (Stadio 5, chiusura della pipeline):} $N^*=3$ per il
correlatore, invariato per $p_\text{readout}\in\{0.00, 0.05, 0.20\}$
--- verificato numericamente, non solo atteso dall'argomento analitico.
La pipeline di rumore sul trimero ad anello è chiusa: cinque stadi,
tutti verificati, inclusa una correzione di percorso (griglia $N$ troppo
grezza allo Stadio 3, corretta retroattivamente --- vedi
\texttt{risultati\_\allowbreak scan\_\allowbreak parametri\_rumore\_trimero\_anello.tex},
Sez.\ ``Correzione'').
\end{quote}

\begin{quote}
\textit{Ultimo passo (addendum su stile e struttura, aggiornato):}
prodotti e verificati end-to-end anche i due pacchetti riproducibili ---
\texttt{pacchetto\_trimero\_\allowbreak riproducibile.zip} (Parte 1) e
\texttt{pacchetto\_trimero\_\allowbreak rumoroso\_\allowbreak riproducibile.zip}
(Parte 2, questa pipeline). Entrambi aggiornati successivamente alla
prima consegna, per coprire contenuti inizialmente assenti: il pacchetto
di Parte 2 include ora anche il readout a shot finiti (non solo la
formula analitica) e il codice per le due estensioni facoltative
(Sez.~\ref{sec:estensioni}); il pacchetto di Parte 1 include ora anche
il modulo RBS (mancava, solo $W$-2q.6 era isolato in un modulo), la
verifica sistematica del termine DM, lo scan completo delle 81
combinazioni correlatore, e il confronto preparazione VQE/esatta ---
tutti in precedenza solo nei notebook esplorativi o non riprodotti.
Ogni script verificato individualmente da uno stato pulito (\texttt{dati/}
e \texttt{figure/} svuotate); numeri rigenerati coincidenti con quelli
riportati in questo documento e nei documenti di stadio.
\end{quote}

\end{document}
